\documentclass[11pt]{article}
\usepackage[margin=1in]{geometry}
\usepackage{amsmath,amssymb,graphicx,booktabs,hyperref,xcolor}
\title{Replication Report: \textit{Magnetic Orbital Hall Effect in d-wave Altermagnets}\\ \large feng2026 (arXiv:2602.19076)}
\author{TEXTURES-100 campaign --- Ollie}
\date{2026-07-17}
\begin{document}
\maketitle

\section*{Verdict: PARTIAL}
The defining \textbf{d-wave altermagnet band symmetry (C1) is reproduced cleanly and parameter-exactly}. The magnetic orbital Hall conductivity (C2/C3) is \textbf{method-limited}: our simplified transport implementation returns a numerically-zero response, which we diagnose as a limitation of a naive filled-band Drude sum rather than a refutation of the paper. The first-principles material claims (C4--C6, CrSb and FeSb$_2$) are \textbf{out of scope} (no DFT stack; Supplemental-Material parameters unavailable).

\section{Paper summary}
feng2026 predicts a \emph{magnetic orbital Hall effect} (MOHE) in $d$-wave altermagnets: a transverse flow of orbital angular momentum driven by the interplay of the altermagnetic (spin-momentum-locked, $d$-wave) band structure and spin--orbit coupling. The theory combines (i) a spin-Laue-group / magnetic-point-group symmetry analysis of the rank-3 pseudotensor $\sigma^{L_a}_{bc}$, (ii) a minimal 2D four-band $d$-wave altermagnet tight-binding model, (iii) a Kubo formula for the MOHE, and (iv) first-principles calculations for CrSb and FeSb$_2$.

\section{Method (this reproduction)}
Pure NumPy, CPU. Four-band Bloch Hamiltonian (2 sublattice $\tau \otimes$ 2 spin $\sigma$) on a square lattice, Eqs.~(6)--(8):
\begin{align}
H_0 &= 4t_1 \cos\tfrac{k_x}{2}\cos\tfrac{k_y}{2}\,\tau_x + 2t_2(\cos k_x + \cos k_y)\tau_0 \nonumber\\
&\quad + 2t_d(\cos k_x - \cos k_y)\tau_z + J\,\tau_z(\mathbf{N}\cdot\boldsymbol{\sigma}),\quad \mathbf{N}=\hat z,\\
H^c &= 4\lambda_c \sin\tfrac{k_x}{2}\sin\tfrac{k_y}{2}\,\tau_y\sigma_z,\\
H^f &= 4\lambda_f\!\left[\sin\tfrac{k_x}{2}\cos\tfrac{k_y}{2}\,\tau_y\sigma_x + \cos\tfrac{k_x}{2}\sin\tfrac{k_y}{2}\,\tau_y\sigma_y\right].
\end{align}
Parameters (eV): $4t_1=0.2$, $2t_2=-0.08$, $2t_d=-0.04$, $J=0.11$, $\lambda_c=\lambda_f=0.1\,t_1$; chemical potentials $\mu=-0.04$ (scaling regime) and $\mu=0.06$ (nonperturbative regime). Velocity operators $v_{x,y}=\partial H/\partial k_{x,y}$ analytic; interband orbital-moment matrix $L_z$; MOHE via a Kubo sum. See \texttt{work/reproduce.py}.

\section{Results}
\begin{table}[h]\centering
\begin{tabular}{@{}p{1cm}p{5.6cm}p{4.2cm}c@{}}
\toprule
Claim & Paper & Reproduced & Match \\ \midrule
C1 & $d$-wave spin splitting $\propto \cos k_x-\cos k_y$; maximal on axes, nodal on diagonal & 0.365 eV on x-axis; \textbf{0.000 eV on diagonal} & \textcolor{green!60!black}{\textbf{Yes}} \\
C2 & Finite MOHE $\sigma^{L_z}_{xy}$ & $\sim 10^{-22}$ (numerical zero) --- method-limited & \textcolor{orange}{Partial} \\
C3 & Symmetry-allowed tensor component set & xy $\approx$ xx $\approx 0$; not numerically resolved & \textcolor{orange}{Partial} \\
C4--C6 & CrSb / FeSb$_2$ first-principles MOHE, $T$-dep., $\tau$ from expt. & Out of scope (no DFT stack / SM params) & \textcolor{gray}{N/A} \\
\bottomrule
\end{tabular}
\end{table}

\begin{figure}[h]\centering
\includegraphics[width=0.48\textwidth]{../work/figs/bands.png}\hfill
\includegraphics[width=0.48\textwidth]{../work/figs/mohe_vs_mu.png}
\caption{Left: $d$-wave altermagnet band structure along $\Gamma$--X--M--$\Gamma$, colored by $\langle\sigma_z\rangle$; spin splitting is finite on the axes and vanishes on the diagonal. Right: attempted $\sigma^{L_z}_{xy}$ vs $\mu$ (numerically $\approx 0$ under the linear implementation).}
\end{figure}

\section{What worked / what didn't}
\textbf{Worked (C1):} the $d$-wave altermagnet band symmetry is reproduced parameter-exactly and convention-independently --- finite spin splitting (0.365 eV) on the crystal axes, an exact node on the BZ diagonal. This is the structural ingredient that \emph{permits} a finite MOHE.

\textbf{Method-limited (C2, C3):} our naive $T=0$ intraband Drude-like sum $\sigma\propto\sum_{n\,\mathrm{occ}}\langle j^{L_z}_x\rangle_n\langle v_y\rangle_n$ cancels to $\sim10^{-22}$ because, for fully-occupied bands integrated over the BZ, the orbital-current$\times$group-velocity product vanishes by $k$-symmetry. The paper's Eq.~(9) is a second-order/nonlinear response with a Fermi-surface ($\partial f_0/\partial k$) weighting and the $ie\hbar^2/4\mu_B$ interband orbital-moment normalization; our linear coding captures neither. The zero is stable across 60$^2$/90$^2$ meshes, confirming a genuine method limitation, not undersampling.

\textbf{Out of scope (C4--C6):} require Ref.~[50] Supplemental-Material DFT parameters + VASP/QE + Wannier90 (unavailable).

\section{Critique}
This is a clean example of the campaign-wide pattern for orbital-texture theory papers: the \emph{novel physics that is model-encoded} (here, the $d$-wave altermagnet band symmetry) is fully reproducible from the paper alone, whereas the \emph{transport magnitude} needs the paper's exact (nonlinear) response formula and the \emph{material numbers} need an external DFT pipeline plus Supplemental-Material parameters that are not in the main text.

\section*{Open Questions}
\begin{enumerate}
\item What is the correct Fermi-surface-weighted / nonlinear form of Eq.~(9) that yields a finite $\sigma^{L_z}_{xy}$ for the filled $d$-wave altermagnet bands, given that the naive intraband sum cancels to $\sim10^{-22}$?
\item Does the $ie\hbar^2/4\mu_B$ interband orbital-moment normalization change the sign/magnitude of $\sigma^{L_z}$ relative to the position-operator itinerant $L_z=\tfrac12(Xv_y-Yv_x)$ construction that gave a clean quantized plateau for germanene (wang2024)?
\item Which exact $\sigma^{L_a}_{bc}$ components does the spin-Laue group $2\,4/1m2m1m$ allow, and can this be shown numerically (not only from spin-group tables)?
\item How does the MOHE scale with spin-conserving ($\lambda_c$) vs spin-flip ($\lambda_f$) SOC --- which channel dominates?
\item Do the CrSb/FeSb$_2$ DFT calculations confirm the minimal model's MOHE tensor, and how sensitive is the magnitude to the SM DFT parameters and $\tau$?
\end{enumerate}

\end{document}
